\documentclass[10pt]{article} 
\renewcommand{\thempfootnote}{\arabic{mpfootnote}}

\usepackage[hidelinks]{hyperref}

\usepackage[utf8]{inputenc}

\usepackage[T1]{fontenc}
\usepackage[ttdefault=true]{AnonymousPro}

\usepackage{multicol}
\setlength{\columnsep}{1cm}

\usepackage{abc}
\usepackage{graphicx}
\usepackage{musicography}

\usepackage[margin=0.5in]{geometry} 
\geometry{letterpaper}

\usepackage{fancyhdr} 
\pagestyle{fancy} % options: empty , plain , fancy
\renewcommand{\headrulewidth}{0pt} % customise the layout...
\lhead{}\chead{}\rhead{}
\lfoot{}\cfoot{\thepage}\rfoot{}

\begin{document}

\thispagestyle{empty}

\begin{abc}[name=gettingupstairs,options={-O= +c -F single.fmt}]
X:1
T:Such a Getting Upstairs
M:2/2
L:1/4
K:G
"G"Bc dg | Bc d-d/d/ | "D"ed cB | AB c d/c/ | 
"G"Bc dg | B/A/B/c/ dd | "D"e-e/d/ cB | AB cd/c/ || 
"G"B/G/A/G/ DG | "D"F/G/A/B/ c/e/d/c/ | "G"B/G/A/G/ DG | "D"F/G/A/F/ "G"Gd/c/ | 
"G"B/G/A/G/ DG | "D"F/G/A/B/ c/e/d/c/ | "G"B/G/A/G/ DG | "D"F/G/A "G"G2 |]
\end{abc}
% \vskip 0.25in
\noindent\makebox[\textwidth][c]{%
  \begin{minipage}{7.0in}
    \setlength{\parindent}{1.5em}

Such a Getting Upstairs is believed to have been composed by Joe
Blackburn in the early 1830s for blackface
performers.\footnote{\url{https://tunearch.org/wiki/Annotation:Getting\_Upstairs\_(1)}}
The song then traveled to England with the American minstrel shows
where it became attached to the English Morris Dance community between
1836 (when it is thought to have been first performed in London) and
the early 20th
century.\footnote{\url{http://pancocojams.blogspot.com/2013/10/such-getting-upstairs-history-of-song.html}
  comment by Billy Weeks on Mudcat, September 25, 2010, 10:46 am} Such
a Getting Upstairs is thought to refer to sexual activities that
occurred on the floor above a tavern and then evolved into meaning
`commotion', `disturbance', or `ruckus' that was notable, exciting,
and
fun.\footnote{\textit{ibid.}
  Mudcat: Getting Upstairs thread general consensus}

Another theory for the origin of `Such a Getting Upstairs' is from a
comical flight upstairs caused by a flood in a small British town as
reported in the \textbf{Sherborne Journal} (Dorset, England) October
23, 1839 and reprinted in \textbf{The Times} (London) October 26,
1839, p.6.\footnote{Unfortunately, I could not verify this article in
  the The Times or the earlier Sherborne Journal since they are behind
  a paywall. I believe the flood became attached to the tune because
  the tune was already popular, rather than the flood inspired the
  tune’s lyrics, but this may be more because of the personal bias
  that sexual relations are more worthy of song than floods (or just
  more fun).}  `Sich a Gittin Up Stars' appears in George P. Knauff's
\textbf{Virginia Reels}, Vol 3 (1844, Baltimore) and `Such a Getting
Up Stairs' appears in Elias Howe's \textbf{Various Collections}
(1840-1861, Boston) firmly placing it in the American fiddling
repertoire during the most active time of the Oregon Trail
(c. 1846-1869).\footnote{Goertzen, Chris. \textbf{Virginia Reels and
    the History of American Fiddling}, p 73.}

The Morris Dance lyrics collected from Headington, Oxfordshire, Cotswolds:

{\parindent0pt
  \raggedright
  ~\\
  Some likes coffee, \\
  Some likes tea, \\
  Some likes a pretty girl, just like me; \\
  Such a getting upstairs and a playing on the fiddle, \\
  Such a getting upstairs I never did see. \\
  ~\\
}

Cecil Sharp included Such a Getting Upstairs in his collection,
English Folk Songs from the Southern Appalachians, as Some Love
Coffee, No 274 (sung by Mrs. Laurel Jones, Burnsville, NC, September
17th, 1918):

{\parindent0pt
  \raggedright
  ~\\
  Some loves coffee, some loves tea, \\
  Some loves money, but they don't love me. \\
  Singing in the lonesome cowboyee, \\
  Singing in the lonesome sea. \\
  ~\\
}

Cecil Sharp's `Some Love Coffee' is in the play-party section of his
collection. This `I love coffee, I love tea' play-party song was a
playground jump roping rhyme pre 1970s; post 1970s as a playground
hand clapping rhyme in the same family as `Shimmy Shimmy Co Co Pa' and
`Down Down Baby'.\footnote{pancojams, \textit{op. cit.} Comment by Azizi Powell, October 29, 2013, 9:11 pm. Yes, I very much remember doing hand clapping rhymes on the school playground with the lyrics `I love coffee, I love tea', `Shimmy Shimmy Co Co Pa', and `Down Down Baby' in Massachusetts in the 1980s.}

  \end{minipage}
  \vspace*{\fill}
}

\end{document}
